%-------------------------
% Resume in Latex
% Author : Sidratul Muntaha Ahmed
% License : MIT
%------------------------

\documentclass[letterpaper,11pt]{article}

\usepackage{latexsym}
\usepackage[empty]{fullpage}
\usepackage{titlesec}
\usepackage{marvosym}
\usepackage[usenames,dvipsnames]{color}
\usepackage{verbatim}
\usepackage{enumitem}
\usepackage[hidelinks]{hyperref}
\usepackage{fancyhdr}
\usepackage[english]{babel}
\usepackage{tabularx}
\usepackage{multicol}
\input{glyphtounicode}


%----------FONT OPTIONS----------
% sans-serif
% \usepackage[sfdefault]{FiraSans}
% \usepackage[sfdefault]{roboto}
% \usepackage[sfdefault]{noto-sans}
% \usepackage[default]{sourcesanspro}

% serif
% \usepackage{CormorantGaramond}
% \usepackage{charter}


\pagestyle{fancy}
\fancyhf{} % clear all header and footer fields
\fancyfoot{}
\renewcommand{\headrulewidth}{0pt}
\renewcommand{\footrulewidth}{0pt}

% Adjust margins
\addtolength{\oddsidemargin}{-0.5in}
\addtolength{\evensidemargin}{-0.5in}
\addtolength{\textwidth}{1in}
\addtolength{\topmargin}{-.5in}
\addtolength{\textheight}{1.15in}

\urlstyle{same}

\raggedbottom
\raggedright
\setlength{\tabcolsep}{0in}

% Sections formatting
\titleformat{\section}{
  \vspace{-4pt}\scshape\raggedright\large
}{}{0em}{}[\color{black}\titlerule \vspace{-5pt}]

% Ensure that generate pdf is machine readable/ATS parsable
\pdfgentounicode=1

%-------------------------
% Custom commands
\newcommand{\resumeItem}[1]{
  \item\small{
    {#1 \vspace{-2pt}}
  }
}

\newcommand{\cveSubheading}[2]{
  \vspace{0pt}\item
    \begin{tabular*}{0.97\textwidth}[t]{l@{\extracolsep{\fill}}r}
      #1 & #2 \\
    \end{tabular*}\vspace{-7pt}
}

\newcommand{\pubSubheading}[3]{
  \vspace{0pt}\item
    \begin{tabular*}{0.97\textwidth}[t]{l@{\extracolsep{\fill}}r}
      \vspace{3pt}
      \textbf{#1} \\ \vspace{3pt}
      #2 \\
      \textit{\small#3}
    \end{tabular*}\vspace{-7pt}
}

\newcommand{\resumeSubheading}[5]{
  \vspace{-2pt}\item
    \begin{tabular*}{0.97\textwidth}[t]{l@{\extracolsep{\fill}}r}
      \textbf{#1} & #2 \\ \vspace{5pt}
      \textit{\small#3} & \textit{\small #4} \\ 
      #5 \\ \vspace{4pt}
    \end{tabular*}\vspace{-7pt}
}

\newcommand{\resumeSubSubheading}[2]{
    \item
    \begin{tabular*}{0.97\textwidth}{l@{\extracolsep{\fill}}r}
      \textit{\small#1} & \textit{\small #2} \\
    \end{tabular*}\vspace{-7pt}
}

\newcommand{\resumeProjectHeading}[3]{
    \item
    \begin{tabular*}{0.97\textwidth}{l@{\extracolsep{\fill}}r}
      \textbf{#1} & #2 \\
      \small#3
    \end{tabular*}\vspace{-7pt}
}

\newcommand{\resumeSubItem}[1]{\resumeItem{#1}\vspace{-4pt}}

\renewcommand\labelitemii{$\vcenter{\hbox{\tiny$\bullet$}}$}

\newcommand{\resumeSubHeadingListStart}{\begin{itemize}[leftmargin=0.15in, label={}, itemsep=6pt, parsep=0pt, topsep=2pt]}
\newcommand{\resumeSubHeadingListEnd}{\end{itemize}}
\newcommand{\resumeItemListStart}{\begin{itemize}[itemsep=6pt, parsep=0pt, topsep=14pt]}
\newcommand{\resumeItemListEnd}{\end{itemize}\vspace{-5pt}}

%-------------------------------------------
%%%%%%  RESUME STARTS HERE  %%%%%%%%%%%%%%%%


\begin{document}

%----------HEADING----------


\begin{center}
    \textbf{\Huge \scshape Cheolwoo Myung} \textbf{ / Staff Engineer, Secure Systems} \\ \vspace{6pt}
    System LSI, Device Solutions \\
    Samsung Electronics \\
    South Korea \\ \vspace{6pt}

    \small Phone: (+82) 10-5500-9036 $|$ Mail: \href{mailto:cheolwmyung@gmail.com}{\underline{cheolwmyung@gmail.com}} $|$
    Company: \href{https://semiconductor.samsung.com}{\underline{Samsung Electronics}}
\end{center}

\vspace{4pt}
%------RESEARCH INTERESTS-------
\section{Summary}

Staff Engineer at Samsung Electronics leading secure element (SE) product
engineering --- supporting Android system security services (e.g., KeyMint
with RKP) and driving PoC engagements to bring the SE total solution
to Android OEMs. Also drives hardware validation techniques for
next-generation SE platforms, spanning a QEMU-based ARM SoC emulator that resolved
software/hardware ambiguity for Samsung Knox Vault and a hardware-in-the-loop
framework built on real chip debug/JTAG interfaces. Brings a Ph.D. research
foundation in Trusted Execution Environments (Intel SGX, AMD SEV-SNP) and
system virtualization security, with fuzzing research on hypervisors and OS
kernels published at USENIX Security and NDSS. Eager to bring this foundation
to the design of next-generation capability-based secure OS architectures.

\vspace{2pt}

%-----------COMPANY PROJECTS-----------
\section{Experience --- Samsung Electronics}
    \resumeSubHeadingListStart
      \resumeProjectHeading
          {\textbf{Secure Element Product Engineering}}
          {Feb. 2026 -- Current}
          {Samsung Electronics}
          \resumeItemListStart
            \resumeItem{Implementing Android KeyMint (StrongBox) with RKP on Samsung's SE secure OS (S3SSE1A/S3SSE2A)}
            \resumeItem{Achieved Android CTS/VTS certification for StrongBox-KeyMint and RKP}
            \resumeItem{Leading PoC engagement with an Android OEM evaluating the SE total solution}
            \resumeItem{Preparing demo applications to showcase the SE solution for future business development}
            \resumeItem{Technologies: Android KeyMint, StrongBox, RKP, Trace32}
          \resumeItemListEnd

      \resumeProjectHeading
          {\textbf{Hardware-in-the-Loop (HIL) Framework for Secure Element Validation}}
          {May 2026 -- Current}
          {Samsung Electronics}
          \resumeItemListStart
            \resumeItem{Designing a hardware-in-the-loop framework using real chip debug/JTAG interfaces}
            \resumeItem{Introduced this HIL validation methodology, established in academic research, as the team's first implementation}
            \resumeItem{Technologies: QEMU, JTAG debugging interface}
          \resumeItemListEnd

      \resumeProjectHeading
          {\textbf{ARM SoC Emulator for Knox Vault Secure OS}}
          {Mar. 2025 -- Jan. 2026}
          {Samsung Electronics}
          \resumeItemListStart
            \resumeItem{Designed a high-fidelity ARM SoC emulator for pre-silicon OS development on Samsung Knox Vault}
            \resumeItem{Resolved approximately 5--6 ambiguous SW/HW issues, driving a hardware workaround in a chip revision}
            \resumeItem{Technologies: QEMU, ARM}
          \resumeItemListEnd
    \resumeSubHeadingListEnd

%-----------PH.D. RESEARCH PROJECTS-----------
\section{Research Projects --- Seoul National University}
    \resumeSubHeadingListStart
      \resumeProjectHeading
          {\textbf{Designing Secure and Efficient TEE Memory Sharing Framework}}
          {Oct. 2024 -- Feb. 2025}
          {Samsung Electronics (Industry Collaboration during Ph.D.)}
          \resumeItemListStart
            \resumeItem{Designed a trusted hypervisor-based framework to securely share TEE memory across a distributed system}
            \resumeItem{Technologies: Intel SGX, KVM, VFIO, RDMA}
          \resumeItemListEnd

      \resumeProjectHeading
          {\textbf{Designing of Android Bootloader Vulnerability Analysis}}
          {Jan. 2023 -- Apr. 2024}
          {Samsung Electronics (Industry Collaboration during Ph.D.)}

          \resumeItemListStart
            \resumeItem{Designed a QEMU/AFL-based fuzzer to detect Samsung Android bootloader vulnerabilities}
          \resumeItemListEnd

      \resumeProjectHeading
          {\textbf{Implementing Secure Model Training Framework}}
          {Mar. 2022 -- Apr. 2023}
          {Seoul National University}
          \resumeItemListStart
            \resumeItem{Implemented a secure model training framework (KVM, VFIO, SEV-SNP) to protect sensitive training data}
          \resumeItemListEnd

      \resumeProjectHeading
          {\textbf{Building Hypervisor Para-Virtualization Driver Fuzzing Framework}}
          {Apr. 2021 -- Oct. 2021}
          {National Security Research Institute of Korea (NSR)}

          \resumeItemListStart
            \resumeItem{Designed a fuzzer to detect vulnerabilities in Hyper-V's para-virtualized (PV) back-end device driver}
          \resumeItemListEnd

      \resumeProjectHeading
          {\textbf{Building Hypervisor Fuzzing Framework}}
          {Apr. 2020 -- Oct. 2020}
          {National Security Research Institute of Korea (NSR)}
          \resumeItemListStart
            \resumeItem{Designed a fuzzer to find vulnerabilities in virtual devices provided by the QEMU hypervisor}
          \resumeItemListEnd

      \resumeProjectHeading
          {\textbf{Developing Data Race Fuzzer for Binary OS Kernel}}
          {Apr. 2019 -- Oct. 2019}
          {National Security Research Institute of Korea (NSR)}
          \resumeItemListStart
            \resumeItem{Designed a QEMU-based fuzzer to detect data race bugs in the Windows 7/10 OS kernel}
          \resumeItemListEnd

    \resumeSubHeadingListEnd

%
%-----------EXPERIENCE-----------
\section{Publications}
  \resumeSubHeadingListStart
    \pubSubheading
      {S3NTRY: A Protection Framework for Secure and Efficient TEE Memory Disaggregated Systems}
      {\textbf{Cheolwoo Myung}, Jaewon Hur, Adil Ahmad, and Byoungyoung Lee}
      {Revising for resubmission}
      \vspace{5pt}

    \pubSubheading
      {DLBox: New Model Training Framework for Protecting Training Data}
      {Jaewon Hur, Juheon Yi, \textbf{Cheolwoo Myung}, Sangyun Kim, Youngki Lee, and Byoungyoung Lee}
      {Network and Distributed System Security Symposium (\textbf{NDSS}), Febuary 2025.}
      \vspace{5pt}

    \pubSubheading
      {MundoFuzz: Hypervisor Fuzzing with Statistical Coverage Testing and Grammar Inference}
      {\textbf{Cheolwoo Myung}, Gwangmu Lee, and Byoungyoung Lee}
      {31st USENIX Security Symposium (\textbf{Security}), August 2022.}
  \resumeSubHeadingListEnd

%
%-----------TECHNICAL SKILLS-----------
\section{Technical Skills}
 \begin{itemize}[leftmargin=0.15in, label={}]
    {\item{
     \textbf{Languages:}{: C, C++, Python, Assembly (ARM, x86)} \\
     \textbf{Frameworks}{: AFL, Syzkaller, QEMU, KVM, Docker, Trace32} \\
     \textbf{Secure Systems}{: Android KeyMint/StrongBox, Remote Key Provisioning (RKP), Secure Element (SE), Intel SGX, AMD SEV-SNP, JTAG/Chip Debug} \\
    }}
\end{itemize}

%-----------EDUCATION-----------
\section{Education}
  \resumeSubHeadingListStart
    \resumeSubheading
        {Seoul National University}{Sep. 2018 - Feb. 2025}
        {Seoul, South Korea}{}
        {Ph.D. in Electrical and Computer Engineering (Advisor: \href{https://lifeasageek.github.io/}{\underline{Byoungyoug Lee}})}
        \vspace{-1em}

    \resumeSubheading
      {Handong Global University}{Mar. 2012 - Aug. 2018}
      {Pohang, South Korea}{}
      {B.S. in Mechanical and Control Engineering}
  \resumeSubHeadingListEnd

%-------------------------------------------
\end{document}
